\documentclass[10pt,a4paper]{amsart}
\usepackage{amsmath,amssymb,amsthm,graphicx}
\usepackage{verbatim}
\newenvironment{reference}{\comment}{\endcomment}
%\newenvironment{reference}{}{}
\newenvironment{slogan}{\comment}{\endcomment}
\newenvironment{history}{\comment}{\endcomment}

% For commutative diagrams we use Xy-pic
\usepackage[all]{xy}

% We use 2cell for 2-commutative diagrams.
\xyoption{2cell}
% Theorem environments.
%
\theoremstyle{plain}
\newtheorem{theorem}[subsection]{Theorem}
\newtheorem{proposition}[subsection]{Proposition}
\newtheorem{lemma}[subsection]{Lemma}

\theoremstyle{definition}
\newtheorem{definition}[subsection]{Definition}
\newtheorem{example}[subsection]{Example}
\newtheorem{exercise}[subsection]{Exercise}
\newtheorem{situation}[subsection]{Situation}

\theoremstyle{remark}
\newtheorem{remark}[subsection]{Remark}
\newtheorem{remarks}[subsection]{Remarks}

\numberwithin{equation}{subsection}

% Macros
%
\def\lim{\mathop{\mathrm{lim}}\nolimits}
\def\colim{\mathop{\mathrm{colim}}\nolimits}
\def\Spec{\mathop{\mathrm{Spec}}}
\def\Hom{\mathop{\mathrm{Hom}}\nolimits}
\def\Ext{\mathop{\mathrm{Ext}}\nolimits}
\def\SheafHom{\mathop{\mathcal{H}\!\mathit{om}}\nolimits}
\def\SheafExt{\mathop{\mathcal{E}\!\mathit{xt}}\nolimits}
\def\Sch{\mathit{Sch}}
\def\Mor{\mathop{\mathrm{Mor}}\nolimits}
\def\Ob{\mathop{\mathrm{Ob}}\nolimits}
\def\Sh{\mathop{\mathit{Sh}}\nolimits}
\def\NL{\mathop{N\!L}\nolimits}
\def\CH{\mathop{\mathrm{CH}}\nolimits}
\def\proetale{{pro\text{-}\acute{e}tale}}
\def\etale{{\acute{e}tale}}
\def\QCoh{\mathit{QCoh}}
\def\Ker{\mathop{\mathrm{Ker}}}
\def\Im{\mathop{\mathrm{Im}}}
\def\Coker{\mathop{\mathrm{Coker}}}
\def\Coim{\mathop{\mathrm{Coim}}}

% Boxtimes
%
\DeclareMathSymbol{\boxtimes}{\mathbin}{AMSa}{"02}

%
% Macros for moduli stacks/spaces
%
\def\QCohstack{\mathcal{QC}\!\mathit{oh}}
\def\Cohstack{\mathcal{C}\!\mathit{oh}}
\def\Spacesstack{\mathcal{S}\!\mathit{paces}}
\def\Quotfunctor{\mathrm{Quot}}
\def\Hilbfunctor{\mathrm{Hilb}}
\def\Curvesstack{\mathcal{C}\!\mathit{urves}}
\def\Polarizedstack{\mathcal{P}\!\mathit{olarized}}
\def\Complexesstack{\mathcal{C}\!\mathit{omplexes}}
% \Pic is the operator that assigns to X its picard group, usage \Pic(X)
% \Picardstack_{X/B} denotes the Picard stack of X over B
% \Picardfunctor_{X/B} denotes the Picard functor of X over B
\def\Pic{\mathop{\mathrm{Pic}}\nolimits}
\def\Picardstack{\mathcal{P}\!\mathit{ic}}
\def\Picardfunctor{\mathrm{Pic}}
\def\Deformationcategory{\mathcal{D}\!\mathit{ef}}
\usepackage{cleveref}
\crefname{originallemma}{original lemma}{original lemmas}
\Crefname{originallemma}{Original lemma}{Original lemmas}
\makeatletter
\newlabel{lemma-regular-quotient-regular}{{00NR}{source}{}{sourcepointer-1}{}}
\newlabel{lemma-regular-quotient-regular@cref}{{[originallemma][0][]00NR}{[1][0][]source}{}{}{}}
\makeatother
\crefname{originallemma}{original lemma}{original lemmas}
\Crefname{originallemma}{Original lemma}{Original lemmas}
\makeatletter
\newlabel{lemma-regular-ring-CM}{{00NQ}{source}{}{sourcepointer-2}{}}
\newlabel{lemma-regular-ring-CM@cref}{{[originallemma][0][]00NQ}{[1][0][]source}{}{}{}}
\makeatother
\crefname{originallemma}{original lemma}{original lemmas}
\Crefname{originallemma}{Original lemma}{Original lemmas}
\makeatletter
\newlabel{proposition-CM-module}{{00N6}{source}{}{sourcepointer-3}{}}
\newlabel{proposition-CM-module@cref}{{[originallemma][0][]00N6}{[1][0][]source}{}{}{}}
\makeatother
\title{Being a local complete intersection is well defined for quotients of regular local rings}
\author{The Stacks Project Authors}
\date{stacks-project:algebra.tex:local-snapshot-20261012; GNU FDL 1.2 or later}
\begin{document}
\maketitle
\noindent\textit{Editorial note.} All delivered mathematical text is the author's own, copied from the pinned Stacks Project chapter file algebra.tex; no statement or proof step was rewritten, reordered or removed.\par\medskip
\noindent\textit{Editorial note.} Preamble provenance, as required for this source: the frozen snapshot stages the chapter's own preamble file beside it (preamble.tex, 217 lines, 8116 bytes, the same shared Stacks-wide preamble the published chow and etale-cohomology cards use); the first line of the pinned chapter file is \textbackslash{}input\{preamble\}, so no preamble text lives inside the chapter itself. The card takes that file's lines 8-18 (the verbatim package and the three \textbackslash{}comment-based environments reference/slogan/history at lines 9-12, section 13.2, together with the Xy-pic setup) and its lines 154 to the runtime end of the file (the native theorem declarations and mathematical macros). Not one character was changed, no package was added, no macro was defined, and the excerpt serves only as preamble, never as a span coordinate.\par\medskip
\noindent\textit{Editorial note.} The same-chapter pointers are located at their \textbackslash{}label lines inside the pinned chapter file; the cross-chapter pointers are located in the project's own tag list staged in the snapshot (sources/stacks-chow/tags.txt), at the verbatim lines carrying their canonical tags.\par\medskip
\noindent\textit{Editorial note.} The counted proof refers to three same-chapter statements by label; all are recorded as navigation stubs pointed at their \textbackslash{}label lines, and the arguments consuming them stay inside the counted span.\par\medskip
\noindent\textit{Editorial note.} The counted spans contain no \textbackslash{}cite, so the card carries no bibliography.\par\medskip
\noindent\textit{Allowed standard prerequisites (editorial).}\begin{itemize}
\item Regular local rings are Cohen-Macaulay. Invoked at the start of the counted proof (line 36267). Allows the number of generators to be compared with the codimension via the theory of Cohen-Macaulay modules.
\item The quotient of a regular local ring by a regular element. Used in the induction at the end of the counted proof (lines 36289-36292). Reduces both the dimension and the minimal number of generators by one.
\end{itemize}\medskip
\ifdefined\hypertarget\hypertarget{sourcepointer-1}{}\fi
\noindent\textit{Source reference.} Original lemma 00NR: sources/stacks-algebra/pinned-original-source.tex, source lines 25967--25967. Navigation pointer for \textbackslash{}ref\{lemma-regular-quotient-regular\} used inside the counted span (line 36270); the argument that consumes it is retained in the span.\par\medskip
\ifdefined\hypertarget\hypertarget{sourcepointer-2}{}\fi
\noindent\textit{Source reference.} Original lemma 00NQ: sources/stacks-algebra/pinned-original-source.tex, source lines 25949--25949. Navigation pointer for \textbackslash{}ref\{lemma-regular-ring-CM\} used inside the counted span (line 36267); the argument that consumes it is retained in the span.\par\medskip
\ifdefined\hypertarget\hypertarget{sourcepointer-3}{}\fi
\noindent\textit{Source reference.} Original lemma 00N6: sources/stacks-algebra/pinned-original-source.tex, source lines 25269--25269. Navigation pointer for \textbackslash{}ref\{proposition-CM-module\} used inside the counted span (line 36269); the argument that consumes it is retained in the span.\par\medskip
\begin{lemma}
\label{lemma-ci-well-defined}
Let $A \to B \to C$ be surjective local ring homomorphisms.
Assume $A$ and $B$ are regular local rings. The following are equivalent
\begin{enumerate}
\item $\Ker(A \to C)$ is generated by a regular sequence,
\item $\Ker(A \to C)$ is generated by $\dim(A) - \dim(C)$ elements,
\item $\Ker(B \to C)$ is generated by a regular sequence, and
\item $\Ker(B \to C)$ is generated by $\dim(B) - \dim(C)$ elements.
\end{enumerate}
\end{lemma}

\begin{proof}
A regular local ring is Cohen-Macaulay, see Lemma \ref{lemma-regular-ring-CM}.
Hence the equivalences (1) $\Leftrightarrow$ (2) and
(3) $\Leftrightarrow$ (4), see Proposition \ref{proposition-CM-module}.
By Lemma \ref{lemma-regular-quotient-regular}
the ideal $\Ker(A \to B)$ can be generated
by $\dim(A) - \dim(B)$ elements.
Hence we see that (4) implies (2).

\medskip\noindent
It remains to show that (1) implies (4). We do this by induction on
$\dim(A) - \dim(B)$. The case $\dim(A) - \dim(B) = 0$ is trivial.
Assume $\dim(A) > \dim (B)$.
Write $I = \Ker(A \to C)$ and $J = \Ker(A \to B)$.
Note that $J \subset I$. Our assumption is that the minimal number
of generators of $I$ is $\dim(A) - \dim(C)$.
Let $\mathfrak m \subset A$ be the maximal
ideal. Consider the maps
$$
J/ \mathfrak m J \to  I / \mathfrak m I \to \mathfrak m /\mathfrak m^2
$$
By Lemma \ref{lemma-regular-quotient-regular} and its proof the
composition is injective. Take any element $x \in J$ which is
not zero in $J /\mathfrak mJ$. By the above and Nakayama's lemma
$x$ is an element of a minimal set of generators of $I$.
Hence we may replace $A$ by $A/xA$ and $I$ by $I/xA$ which
decreases both $\dim(A)$ and the minimal number of generators of $I$
by $1$. Thus we win.
\end{proof}

\end{document}
